% !TeX root = ../main.tex
% Fonte: Lezione 11_12, pp. 2--5.

\addtocontents{toc}{\protect\Needspace{9\protect\baselineskip}}
\section{Costruzioni di sottospazi in \texorpdfstring{$\R^5$}{R5}}

\subsection{Dati e richieste dell'esercizio}
\begin{esercizio}
Sia $V\subseteq\R^5$ il sottospazio definito da
\[
  x_1+x_2+x_3+x_4+x_5=0,
\]
e sia $U\subseteq V$ il sottospazio generato da
\begin{align*}
  u_1&=(3,0,1,-2,-2),\\
  u_2&=(1,-2,3,-1,-1),\\
  u_3&=(3,6,-7,-1,-1).
\end{align*}
Si richiede di:\footnote{Esercizio dell'esame del 7 luglio 2025.}
\begin{enumerate}[label=(\alph*)]
  \item trovare una relazione di dipendenza lineare tra i tre generatori e una base di $U$;
  \item trovare le equazioni cartesiane di $U$;
  \item trovare una base di un sottospazio $W\subseteq\R^5$ tale che $U\oplus W=V$;
  \item trovare una base di un sottospazio $Z\subseteq\R^5$ tale che $\dim Z=3$, $\dim(Z\cap U)=1$ e $\dim(Z\cap W)=1$.
\end{enumerate}
\end{esercizio}

\begin{svolgimento}
\textbf{Appartenenza a $V$.} I generatori appartengono effettivamente a $V$: le somme delle loro coordinate sono rispettivamente $3+0+1-2-2=0$, $1-2+3-1-1=0$ e $3+6-7-1-1=0$. Poiché $V$ è un sottospazio, contiene tutte le loro combinazioni, quindi $U\subseteq V$.

\textbf{Dimensione di $V$.} La dimensione di $V$ è $4$. Infatti la sua equazione equivale a $x_5=-x_1-x_2-x_3-x_4$, perciò
\[
  (x_1,x_2,x_3,x_4,x_5)=\sum_{i=1}^4 x_i(e_i-e_5),
\]
dove $e_1,\ldots,e_5$ sono i vettori coordinati di $\R^5$. I quattro vettori $e_i-e_5$ sono anche indipendenti: le prime quattro coordinate di una combinazione nulla sono i rispettivi coefficienti.
\end{svolgimento}

\subsection{Relazione di dipendenza e base}
\begin{svolgimento}
\textbf{Relazione di dipendenza.} Cerchiamo gli scalari per cui $\alpha_1u_1+\alpha_2u_2+\alpha_3u_3=0$. Il sistema per coordinate è
\[
  \begin{cases}
    3\alpha_1+\alpha_2+3\alpha_3=0,\\
    -2\alpha_2+6\alpha_3=0,\\
    \alpha_1+3\alpha_2-7\alpha_3=0,\\
    -2\alpha_1-\alpha_2-\alpha_3=0,\\
    -2\alpha_1-\alpha_2-\alpha_3=0.
  \end{cases}
\]
La seconda equazione dà $\alpha_2=3\alpha_3$; sostituendo nella terza si ottiene $\alpha_1=-2\alpha_3$. Questi valori soddisfano anche la prima, la quarta e la quinta equazione. Tutte le relazioni sono quindi multipli di quella ottenuta ponendo $\alpha_3=1$:
\[
  -2u_1+3u_2+u_3=0,
  \qquad\text{da cui}\qquad u_3=2u_1-3u_2.
\]
\textbf{Base di $U$.} Eliminando $u_3$ restano due generatori di $U$. Per controllare l'indipendenza, se $au_1+bu_2=0$, la seconda coordinata dà $-2b=0$, dunque $b=0$; la terza dà allora $a=0$. Pertanto $(u_1,u_2)$ è una base di $U$ e $\dim U=2$.
\end{svolgimento}

\subsection{Metodo matriciale per conservare le relazioni}
\begin{svolgimento}
Un modo alternativo di trovare la relazione consiste nel disporre i generatori per righe e affiancare la matrice identità. Il blocco destro registra i coefficienti con cui ogni riga si esprime mediante i tre generatori iniziali:
\[
  M_0=\left(\begin{array}{rrrrr|rrr}
    3&0&1&-2&-2&1&0&0\\
    1&-2&3&-1&-1&0&1&0\\
    3&6&-7&-1&-1&0&0&1
  \end{array}\right).
\]
Sostituiamo la seconda riga con $R_1-3R_2$ e la terza con $R_3-R_1$, eseguendo le stesse operazioni su entrambi i blocchi:
\[
  M_1=\left(\begin{array}{rrrrr|rrr}
    3&0&1&-2&-2&1&0&0\\
    0&6&-8&1&1&1&-3&0\\
    0&6&-8&1&1&-1&0&1
  \end{array}\right).
\]
Ora poniamo $R_3\leftarrow R_3-R_2$:
\[
  M_2=\left(\begin{array}{rrrrr|rrr}
    3&0&1&-2&-2&1&0&0\\
    0&6&-8&1&1&1&-3&0\\
    0&0&0&0&0&-2&3&1
  \end{array}\right).
\]
La terza riga del blocco sinistro è il vettore nullo; i numeri $(-2,3,1)$ nel blocco destro ne descrivono l'espressione $-2u_1+3u_2+u_3=0$. Le prime due righe del blocco sinistro sono indipendenti: in una loro combinazione nulla la prima coordinata annulla il coefficiente della prima riga, poi la seconda annulla quello della seconda. Le operazioni hanno conservato il sottospazio generato, e ne confermano quindi la dimensione $2$.
\end{svolgimento}

\subsection{Dalle equazioni parametriche alle cartesiane}
\begin{svolgimento}
\textbf{Necessità.} Un elemento $x\in U$ è della forma $x=\lambda_1u_1+\lambda_2u_2$. Le equazioni parametriche sono dunque
\[
  \begin{cases}
    x_1=3\lambda_1+\lambda_2,\\
    x_2=-2\lambda_2,\\
    x_3=\lambda_1+3\lambda_2,\\
    x_4=-2\lambda_1-\lambda_2,\\
    x_5=-2\lambda_1-\lambda_2.
  \end{cases}
\]
La seconda e la terza forniscono
\[
  \lambda_2=-\frac{x_2}{2},\qquad
  \lambda_1=x_3-3\lambda_2=x_3+\frac32x_2.
\]
Sostituendo nelle altre tre,
\begin{align*}
  x_1&=3\left(x_3+\frac32x_2\right)-\frac12x_2
        =3x_3+4x_2,\\
  x_4&=-2\left(x_3+\frac32x_2\right)+\frac12x_2
        =-2x_3-\frac52x_2,\\
  x_5&=x_4.
\end{align*}
Otteniamo così le equazioni cartesiane di $U$ in $\R^5$:
\[
  \begin{cases}
    x_1-4x_2-3x_3=0,\\
    2x_4+5x_2+4x_3=0,\\
    x_5-x_4=0.
  \end{cases}
\]
\textbf{Sufficienza.} Le condizioni sono anche sufficienti: dati $x_2$ e $x_3$, i valori di $\lambda_1$ e $\lambda_2$ appena trovati ricostruiscono tali due coordinate; le tre equazioni obbligano le coordinate restanti a coincidere con quelle di $\lambda_1u_1+\lambda_2u_2$.
\end{svolgimento}

\subsection{Costruzione di un complementare in \texorpdfstring{$V$}{V}}
\begin{svolgimento}
\textbf{Scelta dei vettori.} Un complementare $W$ deve avere dimensione $\dim V-\dim U=4-2=2$. Scegliamo
\[
  w_1=(1,0,0,0,-1),\qquad w_2=(0,1,0,-1,0).
\]
Entrambi appartengono a $V$, perché la somma delle coordinate è nulla.

\textbf{Indipendenza.} La relazione $au_1+bu_2+cw_1+dw_2=0$ equivale a
\[
  \begin{cases}
    3a+b+c=0,\\
    -2b+d=0,\\
    a+3b=0,\\
    -2a-b-d=0,\\
    -2a-b-c=0.
  \end{cases}
\]
La seconda e la terza danno $d=2b$ e $a=-3b$. La quarta diventa $6b-b-2b=3b=0$, quindi $b=0$ e di conseguenza $a=d=0$. La prima dà infine $c=0$.

\textbf{Conclusione.} I quattro vettori indipendenti appartengono a $V$, che ha dimensione $4$, dunque ne costituiscono una base. Perciò
\[
  W=\langle w_1,w_2\rangle,\qquad V=U\oplus W.
\]
In particolare $(w_1,w_2)$ è la base richiesta: i vettori aggiunti alla base di $U$ generano tutto ciò che serve per completarla a una base di $V$, e l'indipendenza esclude elementi non nulli in $U\cap W$.
\end{svolgimento}

\subsection{Un sottospazio con due intersezioni assegnate}
\begin{svolgimento}
\textbf{Costruzione.} Per ottenere un sottospazio che incontri sia $U$ sia $W$ in una retta, iniziamo scegliendo un vettore non nullo di ciascuno: $z_1=u_1$ e $z_2=w_1$. Come terzo vettore prendiamo $z_3=e_1=(1,0,0,0,0)$, che non appartiene a $V$ perché la somma delle sue coordinate è $1$. Questa scelta permette di controllare le intersezioni con entrambi i sottospazi.\footnote{La sola condizione $z_3\notin U\cup W$ non basterebbe: per esempio $z_3=u_2+w_1$ è esterno sia a $U$ sia a $W$, ma insieme a $u_1,w_1$ genera anche $u_2$, rendendo l'intersezione con $U$ bidimensionale.} Poniamo dunque
\[
  Z=\langle u_1,w_1,e_1\rangle.
\]
\textbf{Dimensione di $Z$.} Questi tre vettori sono indipendenti. Infatti, in $au_1+bw_1+ce_1=0$, la terza coordinata impone $a=0$, la quinta impone allora $b=0$ e la prima impone $c=0$. Quindi $\dim Z=3$ e $(u_1,w_1,e_1)$ è una base di $Z$.

\textbf{Intersezione con $V$.} Se $z=au_1+bw_1+ce_1\in Z$, la somma delle sue coordinate è $c$, poiché $u_1,w_1\in V$. Ne segue che $z\in V$ se e solo se $c=0$, ossia
\[
  Z\cap V=\langle u_1,w_1\rangle.
\]
\passomath{Z\cap U\subseteq\langle u_1\rangle} Se $z\in Z\cap U$, la sua scrittura $z=au_1+bw_1$ è una scomposizione secondo $V=U\oplus W$. L'unicità di tale scomposizione forza $bw_1=0$, quindi $b=0$.

\passomath{\langle u_1\rangle\subseteq Z\cap U} Ogni multiplo di $u_1$ appartiene sia a $Z$ sia a $U$. Si conclude
\[
  Z\cap U=\langle u_1\rangle,\qquad \dim(Z\cap U)=1.
\]
\passomath{Z\cap W\subseteq\langle w_1\rangle} Se $z\in Z\cap W$, si ha ancora $z=au_1+bw_1$. L’unicità della scomposizione in $U\oplus W$ impone $au_1=0$, quindi $a=0$.

\passomath{\langle w_1\rangle\subseteq Z\cap W} Ogni multiplo di $w_1$ appartiene sia a $Z$ sia a $W$. Pertanto
\[
  Z\cap W=\langle w_1\rangle,\qquad \dim(Z\cap W)=1.
\]
\end{svolgimento}

